\documentclass[10pt,a4paper]{amsart}
\usepackage{amsmath,amssymb,amsthm,graphicx}
% For commutative diagrams we use Xy-pic
\usepackage[all]{xy}

% We use 2cell for 2-commutative diagrams.
\xyoption{2cell}
% Theorem environments.
%
\theoremstyle{plain}
\newtheorem{theorem}[subsection]{Theorem}
\newtheorem{proposition}[subsection]{Proposition}
\newtheorem{lemma}[subsection]{Lemma}

\theoremstyle{definition}
\newtheorem{definition}[subsection]{Definition}
\newtheorem{example}[subsection]{Example}
\newtheorem{exercise}[subsection]{Exercise}
\newtheorem{situation}[subsection]{Situation}

\theoremstyle{remark}
\newtheorem{remark}[subsection]{Remark}
\newtheorem{remarks}[subsection]{Remarks}

\numberwithin{equation}{subsection}

% Macros
%
\def\lim{\mathop{\mathrm{lim}}\nolimits}
\def\colim{\mathop{\mathrm{colim}}\nolimits}
\def\Spec{\mathop{\mathrm{Spec}}}
\def\Hom{\mathop{\mathrm{Hom}}\nolimits}
\def\Ext{\mathop{\mathrm{Ext}}\nolimits}
\def\SheafHom{\mathop{\mathcal{H}\!\mathit{om}}\nolimits}
\def\SheafExt{\mathop{\mathcal{E}\!\mathit{xt}}\nolimits}
\def\Sch{\mathit{Sch}}
\def\Mor{\mathop{\mathrm{Mor}}\nolimits}
\def\Ob{\mathop{\mathrm{Ob}}\nolimits}
\def\Sh{\mathop{\mathit{Sh}}\nolimits}
\def\NL{\mathop{N\!L}\nolimits}
\def\CH{\mathop{\mathrm{CH}}\nolimits}
\def\proetale{{pro\text{-}\acute{e}tale}}
\def\etale{{\acute{e}tale}}
\def\QCoh{\mathit{QCoh}}
\def\Ker{\mathop{\mathrm{Ker}}}
\def\Im{\mathop{\mathrm{Im}}}
\def\Coker{\mathop{\mathrm{Coker}}}
\def\Coim{\mathop{\mathrm{Coim}}}

% Boxtimes
%
\DeclareMathSymbol{\boxtimes}{\mathbin}{AMSa}{"02}

%
% Macros for moduli stacks/spaces
%
\def\QCohstack{\mathcal{QC}\!\mathit{oh}}
\def\Cohstack{\mathcal{C}\!\mathit{oh}}
\def\Spacesstack{\mathcal{S}\!\mathit{paces}}
\def\Quotfunctor{\mathrm{Quot}}
\def\Hilbfunctor{\mathrm{Hilb}}
\def\Curvesstack{\mathcal{C}\!\mathit{urves}}
\def\Polarizedstack{\mathcal{P}\!\mathit{olarized}}
\def\Complexesstack{\mathcal{C}\!\mathit{omplexes}}
% \Pic is the operator that assigns to X its picard group, usage \Pic(X)
% \Picardstack_{X/B} denotes the Picard stack of X over B
% \Picardfunctor_{X/B} denotes the Picard functor of X over B
\def\Pic{\mathop{\mathrm{Pic}}\nolimits}
\def\Picardstack{\mathcal{P}\!\mathit{ic}}
\def\Picardfunctor{\mathrm{Pic}}
\def\Deformationcategory{\mathcal{D}\!\mathit{ef}}
\usepackage{cleveref}
\crefname{originalreference}{original reference}{original references}
\Crefname{originalreference}{Original reference}{Original references}
\makeatletter
\newlabel{lemma-colimit}{{00D5}{source}{}{sourcepointer-1}{}}
\newlabel{lemma-colimit@cref}{{[originalreference][0][]00D5}{[1][0][]source}{}{}{}}
\makeatother
\crefname{originalreference}{original reference}{original references}
\Crefname{originalreference}{Original reference}{Original references}
\makeatletter
\newlabel{lemma-constructible-maps-into-constant-general}{{09Z7}{source}{}{sourcepointer-2}{}}
\newlabel{lemma-constructible-maps-into-constant-general@cref}{{[originalreference][0][]09Z7}{[1][0][]source}{}{}{}}
\makeatother
\crefname{originalreference}{original reference}{original references}
\Crefname{originalreference}{Original reference}{Original references}
\makeatletter
\newlabel{lemma-finite-pushforward-commutes-with-base-change}{{0959}{source}{}{sourcepointer-3}{}}
\newlabel{lemma-finite-pushforward-commutes-with-base-change@cref}{{[originalreference][0][]0959}{[1][0][]source}{}{}{}}
\makeatother
\crefname{originalreference}{original reference}{original references}
\Crefname{originalreference}{Original reference}{Original references}
\makeatletter
\newlabel{lemma-finite-pushforward-constructible}{{095R}{source}{}{sourcepointer-4}{}}
\newlabel{lemma-finite-pushforward-constructible@cref}{{[originalreference][0][]095R}{[1][0][]source}{}{}{}}
\makeatother
\crefname{originalreference}{original reference}{original references}
\Crefname{originalreference}{Original reference}{Original references}
\makeatletter
\newlabel{lemma-integral-cover-trivial-cohomology}{{09ZD}{source}{}{sourcepointer-5}{}}
\newlabel{lemma-integral-cover-trivial-cohomology@cref}{{[originalreference][0][]09ZD}{[1][0][]source}{}{}{}}
\makeatother
\crefname{originalreference}{original reference}{original references}
\Crefname{originalreference}{Original reference}{Original references}
\makeatletter
\newlabel{lemma-torsion-colimit-constructible}{{03SA}{source}{}{sourcepointer-6}{}}
\newlabel{lemma-torsion-colimit-constructible@cref}{{[originalreference][0][]03SA}{[1][0][]source}{}{}{}}
\makeatother
\crefname{originalreference}{original reference}{original references}
\Crefname{originalreference}{Original reference}{Original references}
\makeatletter
\newlabel{proposition-finite-higher-direct-image-zero}{{03QP}{source}{}{sourcepointer-7}{}}
\newlabel{proposition-finite-higher-direct-image-zero@cref}{{[originalreference][0][]03QP}{[1][0][]source}{}{}{}}
\makeatother
\crefname{originalreference}{original reference}{original references}
\Crefname{originalreference}{Original reference}{Original references}
\makeatletter
\newlabel{proposition-leray}{{03QC}{source}{}{sourcepointer-8}{}}
\newlabel{proposition-leray@cref}{{[originalreference][0][]03QC}{[1][0][]source}{}{}{}}
\makeatother
\title{An arbitrary class on a closed subscheme of an affine scheme is killed by passing to a larger torsion sheaf}
\author{The Stacks Project Authors}
\date{stacks-project:etale-cohomology.tex:local-snapshot-20261011; GNU FDL 1.2 or later}
\begin{document}
\maketitle
\noindent\textit{Editorial note.} All delivered mathematical text is the author's own, copied from the pinned Stacks Project chapter file etale-cohomology.tex; no statement or proof step was rewritten, reordered or removed.\par\medskip
\noindent\textit{Editorial note.} The counted proof is the author's complete argument: the proof is written out in full: the reduction to constructible coefficients with the pushout argument, the reduction to a finite pushforward of constants, the Leray computation on the base-changed cover, and the closing stalk computation proving injectivity.\par\medskip
\noindent\textit{Editorial note.} The label lemma-efface-cohomology-on-closed-by-finite-cover names the counted statement, and the proof environment immediately following it in the source is delivered verbatim as its proof.\par\medskip
\noindent\textit{Editorial note.} The counted statement is carried by the author's own lemma environment in the chapter source.\par\medskip
\noindent\textit{Editorial note.} Pre-open gate: python operations/agent-authoring/tools/precard\_gate.py --source etale-cohomology.tex --label lemma-efface-cohomology-on-closed-by-finite-cover --start 14900 --end 14967 returned RESULT=OK (occupancy FREE).\par\medskip
\noindent\textit{Editorial note.} The statement and proof are copied from the pinned chapter file; the cited inputs are kept as external pointers with their canonical Stacks tags.\par\medskip
\noindent\textit{Allowed standard prerequisites (editorial).}\begin{itemize}
\item Torsion sheaves are filtered colimits of constructible ones. Same chapter, the lemmas describing a torsion sheaf through its constructible subsheaves and computing cohomology of a filtered colimit. Supplies the reduction step producing a constructible G with a class mapping to the given xi.
\item Embedding of a constructible sheaf into finite pushforwards of constants. Same chapter, the general embedding lemma for constructible sheaves into finite products or sums of pushforwards of constant sheaves along finite finitely presented morphisms. Supplies the reduction to F = f\_*underline M.
\item Leray and vanishing of higher direct images for finite morphisms; base change of finite pushforwards. Same chapter, the Leray spectral sequence and the vanishing of the higher direct images of a finite morphism, together with the base change compatibility of a finite pushforward. Supplies the identification H\textasciicircum{}q(Z, f\_*M|\_Z) = H\textasciicircum{}q(Y x\_X Z, M).
\item The integral cover killing cohomology of constant coefficients. Same chapter, the lemma producing a finite surjective finitely presented cover over which a given class with finite constant coefficients dies. Supplies the cover Y' -> Y and hence the injective map f\_*M -> f'\_*M.
\end{itemize}\medskip
\ifdefined\hypertarget\hypertarget{sourcepointer-1}{}\fi
\noindent\textit{Source reference.} Original reference 00D5: sources/stacks-etale-cohomology/pinned-original-source.tex, source lines 7427--7427. Prerequisite pointer invoked at proof line(s) 14911; it names a standard fact of the same chapter that the counted argument uses, and every step of the counted proof built on it is retained verbatim in the delivered body.\par\medskip
\ifdefined\hypertarget\hypertarget{sourcepointer-2}{}\fi
\noindent\textit{Source reference.} Original reference 09Z7: sources/stacks-etale-cohomology/pinned-original-source.tex, source lines 13028--13028. Prerequisite pointer invoked at proof line(s) 14928; it names a standard fact of the same chapter that the counted argument uses, and every step of the counted proof built on it is retained verbatim in the delivered body.\par\medskip
\ifdefined\hypertarget\hypertarget{sourcepointer-3}{}\fi
\noindent\textit{Source reference.} Original reference 0959: sources/stacks-etale-cohomology/pinned-original-source.tex, source lines 8292--8292. Prerequisite pointer invoked at proof line(s) 14936; it names a standard fact of the same chapter that the counted argument uses, and every step of the counted proof built on it is retained verbatim in the delivered body.\par\medskip
\ifdefined\hypertarget\hypertarget{sourcepointer-4}{}\fi
\noindent\textit{Source reference.} Original reference 095R: sources/stacks-etale-cohomology/pinned-original-source.tex, source lines 12528--12528. Prerequisite pointer invoked at proof line(s) 14933; it names a standard fact of the same chapter that the counted argument uses, and every step of the counted proof built on it is retained verbatim in the delivered body.\par\medskip
\ifdefined\hypertarget\hypertarget{sourcepointer-5}{}\fi
\noindent\textit{Source reference.} Original reference 09ZD: sources/stacks-etale-cohomology/pinned-original-source.tex, source lines 14467--14467. Prerequisite pointer invoked at proof line(s) 14947; it names a standard fact of the same chapter that the counted argument uses, and every step of the counted proof built on it is retained verbatim in the delivered body.\par\medskip
\ifdefined\hypertarget\hypertarget{sourcepointer-6}{}\fi
\noindent\textit{Source reference.} Original reference 03SA: sources/stacks-etale-cohomology/pinned-original-source.tex, source lines 12216--12216. Prerequisite pointer invoked at proof line(s) 14911; it names a standard fact of the same chapter that the counted argument uses, and every step of the counted proof built on it is retained verbatim in the delivered body.\par\medskip
\ifdefined\hypertarget\hypertarget{sourcepointer-7}{}\fi
\noindent\textit{Source reference.} Original reference 03QP: sources/stacks-etale-cohomology/pinned-original-source.tex, source lines 8245--8245. Prerequisite pointer invoked at proof line(s) 14942; it names a standard fact of the same chapter that the counted argument uses, and every step of the counted proof built on it is retained verbatim in the delivered body.\par\medskip
\ifdefined\hypertarget\hypertarget{sourcepointer-8}{}\fi
\noindent\textit{Source reference.} Original reference 03QC: sources/stacks-etale-cohomology/pinned-original-source.tex, source lines 8176--8176. Prerequisite pointer invoked at proof line(s) 14940; it names a standard fact of the same chapter that the counted argument uses, and every step of the counted proof built on it is retained verbatim in the delivered body.\par\medskip
\begin{lemma}
\label{lemma-efface-cohomology-on-closed-by-finite-cover}
Let $X$ be an affine scheme. Let $\mathcal{F}$ be a torsion abelian sheaf
on  $X_\etale$. Let $Z \subset X$ be a closed subscheme. Let
$\xi \in H^q_\etale(Z, \mathcal{F}|_Z)$ for some $q > 0$.
Then there exists an injective map $\mathcal{F} \to \mathcal{F}'$
of torsion abelian sheaves on $X_\etale$ such that
the image of $\xi$ in $H^q_\etale(Z, \mathcal{F}'|_Z)$ is zero.
\end{lemma}

\begin{proof}
By Lemmas \ref{lemma-torsion-colimit-constructible} and \ref{lemma-colimit}
we can find a map $\mathcal{G} \to \mathcal{F}$ with $\mathcal{G}$
a constructible abelian sheaf and $\xi$ coming from an element $\zeta$ of
$H^q_\etale(Z, \mathcal{G}|_Z)$. Suppose we can find an injective map
$\mathcal{G} \to \mathcal{G}'$ of torsion abelian sheaves on $X_\etale$
such that the image of $\zeta$ in $H^q_\etale(Z, \mathcal{G}'|_Z)$ is zero.
Then we can take $\mathcal{F}'$ to be the pushout
$$
\mathcal{F}' = \mathcal{G}' \amalg_{\mathcal{G}} \mathcal{F}
$$
and we conclude the result of the lemma holds. (Observe that restriction
to $Z$ is exact, so commutes with finite limits and colimits and moreover
it commutes with arbitrary colimits as a left adjoint to pushforward.)
Thus we may assume $\mathcal{F}$ is constructible.

\medskip\noindent
Assume $\mathcal{F}$ is constructible. By
Lemma \ref{lemma-constructible-maps-into-constant-general}
it suffices to prove the result when $\mathcal{F}$
is of the form $f_*\underline{M}$ where $M$ is a finite abelian group
and $f : Y \to X$ is a finite morphism of finite presentation
(such sheaves are still constructible by
Lemma \ref{lemma-finite-pushforward-constructible}
but we won't need this).
Since formation of $f_*$ commutes with any base change
(Lemma \ref{lemma-finite-pushforward-commutes-with-base-change})
we see that the restriction of $f_*\underline{M}$ to $Z$ is
equal to the pushforward of $\underline{M}$ via
$Y \times_X Z \to Z$. By the Leray spectral sequence
(Proposition \ref{proposition-leray})
and vanishing of higher direct images
(Proposition \ref{proposition-finite-higher-direct-image-zero}),
we find
$$
H^q_\etale(Z, f_*\underline{M}|_Z) = H^q_\etale(Y \times_X Z, \underline{M}).
$$
By Lemma \ref{lemma-integral-cover-trivial-cohomology}
we can find a finite surjective morphism $Y' \to Y$ of finite presentation
such that $\xi$ maps to zero in $H^q(Y' \times_X Z, \underline{M})$.
Denoting $f' : Y' \to X$ the composition $Y' \to Y \to X$ we claim
the map
$$
f_*\underline{M} \longrightarrow f'_*\underline{M}
$$
is injective which finishes the proof by what was said above.
To see the desired injectivity we can look at stalks. Namely,
if $\overline{x} : \Spec(k) \to X$ is a geometric point, then
$$
(f_*\underline{M})_{\overline{x}} =
\bigoplus\nolimits_{f(\overline{y}) = \overline{x}} M
$$
by Proposition \ref{proposition-finite-higher-direct-image-zero}
and similarly for the other sheaf.
Since $Y' \to Y$ is surjective and finite we see that
the induced map on geometric points lifting $\overline{x}$ is
surjective too and we conclude.
\end{proof}

\end{document}
